\scanpage{55}

$U\longrightarrow\overline H^{+}$ sont exactement les cartes topologiques (i.e. \deleted{applications continues injectives}, \inserted{plongements \deleted{inj.} \deleted{\unsure{de}} \inserted{ouverts}, ou encore \unsure{continues inj.|cartes au sens top.|continus inj.}} appliquant bord dans bord) qui sont conformes dans $\operatorname{Int}U$. Cela \deleted{\unsure{provient du|résulte du|précédent du}} \inserted{équivaut au}

\uline{Lemme 1} Soit $U$ un ouvert de $\overline H^{+}$, alors \deleted{toute} \deleted{application \inserted{\deleted{\unsure{continue}}} injective} plongement ouvert $f:U\longrightarrow\overline H$ \deleted{fonction continue sur $V$,} qui est holo dans $U\cap H$, est holo (i.e. se prolonge holomorphiquement dans $\mathbb C$, au voisinage de tt pt de $U\cap\partial H=U\cap\mathbb R$).

C'est un cas particulier du \uline{lemme de \unsure{Schwartz|Schwarz}}, qui dit que tte fonction \deleted{holo} continue complexe $f$ dans $U$, qui est \uline{réelle} sur $U\cap\mathbb R$, est holomorphe.

Pour terminer la dém du \deleted{lemme} \inserted{th}, il faut voir encore que le $f$ du lemme est bien (au voisinage de tt pt $x$ de $U\cap\partial H$) une carte holomorphe de $\deleted{\unsure{H|B|\partial}}U$ (i.e. de $H$) en $x$, i.e. que le germe \inserted{en $x$} de fonction holomorphe \uline{sur $\mathbb C$} défini par $f$ est inversible. Mais \deleted{\unsure{il}} ce germe n'est pas constant, \deleted{et la structure} \inserted{de par son} développement
% archive sheet number: 57
